\section{Materiais e métodos}

\pendente{Descrever o que foi efetivamente feito, no passado, e não o que se
pretendia fazer. Divergências em relação ao planejado no TCC I são material
legítimo do texto: registre a mudança e o motivo.}

\subsection{Arquitetura da solução}

\pendente{Diagrama e descrição dos componentes.}

\subsection{Modelagem de dados}

\pendente{Modelo conceitual e lógico, incluindo o mecanismo de trilha de
auditoria que permite reproduzir uma pesquisa de preços na data em que foi
realizada.}

\subsection{Integração com dados públicos de contratações}

\pendente{Endpoints consumidos, parâmetros de consulta, tratamento de falhas e
estratégia de cache.}

\subsection{Tratamento estatístico dos preços}

\pendente{Critério de exclusão de valores discrepantes e cálculo dos
parâmetros, com a justificativa normativa de cada escolha.}

\subsection{Geração assistida de minutas}

\pendente{Como os documentos são produzidos e como sua conformidade é
verificada.}

\subsection{Procedimento de avaliação}

\pendente{Participantes, tarefas, métricas e instrumentos. Detalhar o suficiente
para que outra pessoa possa repetir o experimento.}
